\documentclass[../book.tex]{subfiles}

\begin{document}

    \section{二重积分}

    \subsection{概念、性质与对称性}

    \subsubsection{概念}

    设 $f(x,y)$ 是有界闭区域 $D$ 上的有界函数，将 $D$ 任意{\color{red}{分割}}成 $n$ 个小闭区域 $\Delta\sigma_1, \Delta\sigma_2, \cdots, \Delta\sigma_n$，其中 $\Delta\sigma_i$ 既表示第 $i$ 个小闭区域，也表示它的面积。在每个 $\Delta\sigma_i$ 上任取一点 $(\xi_i, \eta_i)$，作乘积$f(\xi_i, \eta_i) \Delta\sigma_i$（{\color{red}{近似}}），并{\color{red}{作和}}
    \[
        \sum_{i=1}^{n} f(\xi_i, \eta_i) \Delta\sigma_i
    \]
    如果当各小闭区域的直径中的最大值 $\lambda \to 0$ 时，这和的{\color{red}{极限}}总存在（与$\Delta\sigma_i$的分法及点$(\xi_i, \eta_i)$的取法均无关），则称此极限为函数 $f(x,y)$ 在闭区域 $D$ 上的\textbf{二重积分}，记为
    \[
        \iint\limits_D f(x,y) \mathrm{d}\sigma = \lim_{\lambda \to 0} \sum_{i=1}^{n} f(\xi_i, \eta_i) \Delta\sigma_i
    \]
    其中 $f(x,y)$ 叫做\textbf{被积函数}，$f(x,y)\mathrm{d}\sigma$ 叫做\textbf{被积表达式}，$\mathrm{d}\sigma$ 叫做\textbf{面积元素}，$x$ 与 $y$ 叫做\textbf{积分变量}，$D$ 叫做\textbf{积分区域}，$\sum\limits_{i=1}^{n} f(\xi_i, \eta_i) \Delta\sigma_i$叫做\textbf{积分和}

    \begin{enumerate}
        \item 几何意义：如果 $f(x,y) \geq 0$，则二重积分 $\displaystyle\iint\limits_D f(x,y) \mathrm{d}\sigma$ 表示以 $D$ 为底、以曲面 $z = f(x,y)$ 为顶的曲顶柱体的体积
        \item 如果 $f(x, y)$ 是负的，柱体在$xOy$面的下方，二重积分的绝对值仍等于柱体的体积，但二重积分的值是负的
        \item 如果 $f(x, y)$ 在$D$的若干部分区域上是正的，而在其他区域上是负的，那么，$f(x, y)$ 在 $D$ 上的二重积分就等于 $xOy$ 面上方的柱体体积减去 $xOy$ 面下方的柱体体积所得之差
    \end{enumerate}

    \subsubsection{性质}

    \begin{enumerate}
        \item \textbf{求区域面积}：$\displaystyle\iint\limits_D 1 \mathrm{d}\sigma = \iint\limits_D \mathrm{d}\sigma = S_D$（区域 $D$ 的面积）

        \item \textbf{可积函数必有界}：若 $f(x,y)$ 在 $D$ 上可积，则 $f(x,y)$ 在 $D$ 上有界

        \item \textbf{积分的线性性质}：设 $k_1, k_2$ 为常数，则
        \[
            \iint\limits_D [k_1 f(x,y) + k_2 g(x,y)] \mathrm{d}\sigma = k_1 \iint\limits_D f(x,y) \mathrm{d}\sigma + k_2 \iint\limits_D g(x,y) \mathrm{d}\sigma
        \]

        \item \textbf{积分的可加性}：如果闭区域 $D$ 被有限条曲线分为有限个部分闭区域，则在 $D$ 上的二重积分等于在各部分闭区域上的二重积分的和。设 $D = D_1 \cup D_2, D_1 \cap D_2 = \varnothing$，则
        \[
            \iint\limits_D f(x,y) \mathrm{d}\sigma = \iint\limits_{D_1} f(x,y) \mathrm{d}\sigma + \iint\limits_{D_2} f(x,y) \mathrm{d}\sigma
        \]

        \item \textbf{积分的保号性}：
        \begin{circlenum}
            \item 若在 $D$ 上 $f(x,y) \geq 0$，则 $\displaystyle\iint\limits_D f(x,y) \mathrm{d}\sigma \geq 0$
            \item 若在 $D$ 上 $f(x,y) \leq g(x,y)$，则 $\displaystyle\iint\limits_D f(x,y) \mathrm{d}\sigma \leq \iint\limits_D g(x,y) \mathrm{d}\sigma$
            \item $\displaystyle\left| \iint\limits_D f(x,y) \mathrm{d}\sigma \right| \leq \iint\limits_D |f(x,y)| \mathrm{d}\sigma$
        \end{circlenum}

        \item \textbf{二重积分的估值定理}：设 $M$ 和 $m$ 分别是 $f(x,y)$ 在闭区域 $D$ 上的最大值和最小值（多元函数微分学的应用），$S_D$ 是 $D$ 的面积，则
        \[
            mS_D \leq \iint\limits_D f(x,y) \mathrm{d}\sigma \leq MS_D
        \]

        \item $\color{red}{\bigstar}$\textbf{二重积分的中值定理}：设 $f(x,y)$ 在闭区域 $D$ 上连续，$S_D$ 是 $D$ 的面积，则在 $D$ 上至少存在一点 $(\xi, \eta)$，使得
        \[
            \iint\limits_D f(x,y) \mathrm{d}\sigma = f(\xi, \eta) S_D
        \]
    \end{enumerate}

    \subsubsection{对称性}

    \begin{enumerate}
        \item \textbf{普通对称性（偶倍奇零）}：
        \begin{itemize}
            \item 设区域 $D$ 关于 $x = a$ 对称，$(x, y)$ 关于 $x = a$ 的对称点为 $(2a-x, y)$，则
            \[
                \iint\limits_D f(x,y) \mathrm{d}\sigma = \begin{cases}
                                                             2\iint\limits_{D_1} f(x,y) \mathrm{d}\sigma, & f(x,y) = f(2a-x, y) \\
                                                             0, & f(x,y) = -f(2a-x, y)
                \end{cases}
            \]
            其中 $D_1$ 为 $D$ 在 $x \geq a$（或 $x \leq a$）一侧的部分

            \item 设区域 $D$ 关于 $y = a$ 对称，$(x, y)$ 关于 $y = a$ 的对称点为 $(x, 2a - y)$，则
            \[
                \iint\limits_D f(x,y) \mathrm{d}\sigma = \begin{cases}
                                                             2\iint\limits_{D_1} f(x,y) \mathrm{d}\sigma, & f(x,y) = f(x, 2a-y) \\
                                                             0, & f(x,y) = -f(x, 2a-y)
                \end{cases}
            \]
            其中 $D_1$ 为 $D$ 在 $y \geq a$（或 $y \leq a$）一侧的部分

            \item \textbf{设区域 $D$ 关于原点对称}，则
            \[
                \iint\limits_D f(x,y) \mathrm{d}\sigma = \begin{cases}
                                                             2\iint\limits_{D_1} f(x,y) \mathrm{d}\sigma, & f(x,y) = f(-x,-y) \\
                                                             0, & f(x,y) = -f(-x,-y)
                \end{cases}
            \]
            其中 $D_1$ 为 $D$ 在 $x \geq 0$（或 $y \geq 0$）一侧的部分

            \item 设区域 $D$ 关于 $y = x$ 对称，则
            \[
                \iint\limits_D f(x,y) \mathrm{d}\sigma = \begin{cases}
                                                             2\iint\limits_{D_1} f(x,y) \mathrm{d}\sigma, & f(x,y) = f(y, x) \\
                                                             0, & f(x,y) = -f(y,x)
                \end{cases}
            \]
            其中 $D_1$ 是 $D$ 关于 $y \geq x$ 对称的半个部分

        \end{itemize}

        \item \textbf{轮换对称性}：若 $D$ 关于 $y = x$ 对称，则
        \[
            \iint\limits_D f(x,y) \mathrm{d}\sigma = \iint\limits_D f(y,x) \mathrm{d}\sigma
        \]

        \begin{circlenum}
            \item 特别地，$\displaystyle\iint\limits_D f(x) \mathrm{d}\sigma = \iint\limits_D f(y) \mathrm{d}\sigma$
            \item 若 $f(x, y) + f(y, x) = a$，则
            \[
                I = \frac{1}{2} \iint\limits_D [f(x, y) + f(y, x)] \mathrm{d}x\mathrm{d}y = \frac{1}{2} \iint\limits_{D} a \mathrm{d}x\mathrm{d}y = \frac{a}{2}S_D
            \]
        \end{circlenum}

    \end{enumerate}

    \subsection{计算}

    \subsubsection{直角坐标系下的计算方法}

    \begin{enumerate}
        \item $\color{red}{\bigstar}$定限口诀：后积先定限，限内画条线，先交写下限，后交写上限
        \item \textbf{X 型区域}（先 $y$ 后 $x$）：若积分区域 $D$ 可表示为
        \[
            D = \{(x,y) \mid a \leq x \leq b, \varphi_1(x) \leq y \leq \varphi_2(x)\}
        \]
        则
        \[
            \iint\limits_D f(x,y) \mathrm{d}x\mathrm{d}y = \int_a^b \mathrm{d}x \int_{\varphi_1(x)}^{\varphi_2(x)} f(x,y) \mathrm{d}y
        \]

        \item \textbf{Y 型区域}（先 $x$ 后 $y$）：若积分区域 $D$ 可表示为
        \[
            D = \{(x,y) \mid c \leq y \leq d, \psi_1(y) \leq x \leq \psi_2(y)\}
        \]
        则
        \[
            \iint\limits_D f(x,y) \mathrm{d}x\mathrm{d}y = \int_c^d \mathrm{d}y \int_{\psi_1(y)}^{\psi_2(y)} f(x,y) \mathrm{d}x
        \]
        \item $\color{red}{\bigstar}$此处的下限都必须小于等于上限
        \item 计算二重积分的{\color{red}{关键}}是确定积分限，为此，要画好积分区域$D$的边界图形，当$D$的边界图形不易画出时，要写出$D$的不等式表达式，从而确定上下限
    \end{enumerate}

    \subsubsection{极坐标系下的计算方法}

    \paragraph{基本原理}

    极坐标系适用于与中心对称图像有关的积分。因为是中心对称图像，一般先积 $r$，后积 $\theta$

    \paragraph{坐标变换公式}

    \textbf{直角坐标 $\to$ 极坐标}（{\color{red}{换元}}）：
    \[
        \begin{cases}
            x = r\cos\theta, \\
            y = r\sin\theta,
        \end{cases}
        \qquad \gmath{\mathrm{d}\sigma = r\,\mathrm{d}r\,\mathrm{d}\theta}
    \]

    \vspace{0.5em}

    \paragraph{$\color{red}{\bigstar}$ 定限口诀}

    \textbf{后积先定限，限内画条线，先交写下限，后交写上限}

    \vspace{0.5em}

    \paragraph{三种情形}

    \begin{enumerate}
        \item \textbf{极点$O$在区域$D$外}

        若 $D = \{(r,\theta) \mid \textcolor{blue}{\alpha \leq \theta \leq \beta},\, \textcolor{red}{r_1(\theta) \leq r \leq r_2(\theta)}\}$，则
        \[
            \iint\limits_D f(r\cos\theta, r\sin\theta)\, r\,\mathrm{d}r\,\mathrm{d}\theta
            = \int_{\textcolor{blue}{\alpha}}^{\textcolor{blue}{\beta}} \mathrm{d}\theta \int_{\textcolor{red}{r_1(\theta)}}^{\textcolor{red}{r_2(\theta)}} f(r\cos\theta, r\sin\theta)\, r\,\mathrm{d}r
        \]

        \item \textbf{极点$O$在区域$D$边界上}

        若 $D = \{(r,\theta) \mid \textcolor{blue}{\alpha \leq \theta \leq \beta},\, \textcolor{red}{0 \leq r \leq r(\theta)}\}$，则
        \[
            \iint\limits_D f(r\cos\theta, r\sin\theta)\, r\,\mathrm{d}r\,\mathrm{d}\theta
            = \int_{\textcolor{blue}{\alpha}}^{\textcolor{blue}{\beta}} \mathrm{d}\theta \int_{\textcolor{red}{0}}^{\textcolor{red}{r(\theta)}} f(r\cos\theta, r\sin\theta)\, r\,\mathrm{d}r
        \]

        \item \textbf{极点$O$在区域$D$内}

        若 $D = \{(r,\theta) \mid \textcolor{blue}{0 \leq \theta \leq 2\pi},\, \textcolor{red}{0 \leq r \leq r(\theta)}\}$，则
        \[
            \iint\limits_D f(r\cos\theta, r\sin\theta)\, r\,\mathrm{d}r\,\mathrm{d}\theta
            = \int_{\textcolor{blue}{0}}^{\textcolor{blue}{2\pi}} \mathrm{d}\theta \int_{\textcolor{red}{0}}^{\textcolor{red}{r(\theta)}} f(r\cos\theta, r\sin\theta)\, r\,\mathrm{d}r
        \]
    \end{enumerate}

    \subsubsection{极坐标系与直角坐标系选择的一般原则}

    \paragraph{优先选用极坐标的情形}
    \begin{enumerate}
        \item $\color{red}{\bigstar}$被积函数是否为 $f(x^2 + y^2)$、$f(\dfrac{y}{x})$、$f(\dfrac{x}{y})$ 等形式
        \item 积分区域 $D$ 是圆域、圆环、扇形或其部分
        \item 边界曲线用极坐标表示更简单（如 $x^2 + y^2 = a^2$ 变为 $r = a$）
    \end{enumerate}

    \subsubsection{极坐标系与直角坐标系的互相转化}

    \begin{enumerate}
        \item \textbf{直角坐标 $\to$ 极坐标}：
        \[
            \begin{cases}
                x = r\cos\theta, \\
                y = r\sin\theta,
            \end{cases}
            \qquad \gmath{\mathrm{d}\sigma = r\,\mathrm{d}r\,\mathrm{d}\theta}
        \]

        \item \textbf{极坐标 $\to$ 直角坐标}：
        \[
            r = \sqrt{x^2 + y^2}, \quad \theta = \arctan\frac{y}{x}
        \]
        \item 画出区域$D$的边界图形，做好上限、下限的转化
    \end{enumerate}

    \subsubsection{换元法}

    设 $f(x,y)$ 在 $D_{xy}$ 上连续，变换 $\begin{cases}
                                             x = x(u,v) \\ y = y(u,v)
    \end{cases}$ 将 $uOv$ 平面上的闭区域 $D_{uv}$ 变为 $xOy$ 平面上的闭区域 $D_{xy}$，且雅可比行列式
    \[
        J = \frac{\partial(x,y)}{\partial(u,v)} = \begin{vmatrix}
                                                      \dfrac{\partial x}{\partial u} & \dfrac{\partial x}{\partial v} \\[6pt]
                                                      \dfrac{\partial y}{\partial u} & \dfrac{\partial y}{\partial v}
        \end{vmatrix} \neq 0
    \]
    则有
    \[
        \iint\limits_{D_{xy}} f(x,y) \mathrm{d}x\mathrm{d}y = \iint\limits_{D_{uv}} f[x(u,v), y(u,v)] |J| \mathrm{d}u\mathrm{d}v
    \]

    二重积分换元法中的{\color{red}{三换}}:
    \begin{enumerate}
        \item 换积分区域
        \item 换被积函数
        \item 换积分变量
    \end{enumerate}

    \subsection{结论}

    \begin{enumerate}
        \item \textbf{椭圆}

        \begin{circlenum}
            \item \textbf{标准方程}：
            \begin{itemize}
                \item 圆心在原点：$\displaystyle\frac{x^2}{a^2} + \frac{y^2}{b^2} = 1$（$a > b > 0$）
                \item 圆心在 $(h, k)$：$\displaystyle\frac{(x-h)^2}{a^2} + \frac{(y-k)^2}{b^2} = 1$
            \end{itemize}
            其中 $a$ 为长半轴，$b$ 为短半轴

            \item \textbf{面积}：$S = \pi ab$
        \end{circlenum}
        \item 当$x \to 0^+$时，$f(x) \sim g(x)$，则$\displaystyle\int_0^x f(t)\,\mathrm{d}t \sim \int_0^x g(t)\,\mathrm{d}t$
        \item \textbf{高斯积分}（概率积分）

        \begin{circlenum}
            \item \textbf{标准形式}：
            \[
                \int_{-\infty}^{+\infty} e^{-x^2} \mathrm{d}x = \sqrt{\pi}
            \]

            等价形式：
            \[
                \int_{0}^{+\infty} e^{-x^2} \mathrm{d}x = \frac{\sqrt{\pi}}{2}
            \]

            \item $\color{red}{\bigstar}$\textbf{计算方法}（利用二重积分与极坐标）：

            设 $I = \displaystyle\int_{0}^{+\infty} e^{-x^2} \mathrm{d}x$，则
            \[
                I^2 = \int_{0}^{+\infty} e^{-x^2} \mathrm{d}x \cdot \int_{0}^{+\infty} e^{-y^2} \mathrm{d}y = \int_{0}^{+\infty}\int_{0}^{+\infty} e^{-(x^2+y^2)} \mathrm{d}x\mathrm{d}y
            \]

            转化为极坐标，区域为第一象限：$0 \leq \theta \leq \frac{\pi}{2}$，$0 \leq r < +\infty$，则
            \[
                I^2 = \int_{0}^{\frac{\pi}{2}} \mathrm{d}\theta \int_{0}^{+\infty} e^{-r^2} \cdot r\,\mathrm{d}r = \frac{\pi}{2} \cdot \left[-\frac{1}{2}e^{-r^2}\right]_{0}^{+\infty} = \frac{\pi}{4}
            \]

            因此 $I = \dfrac{\sqrt{\pi}}{2}$

            \item \textbf{一般形式}：
            \[
                \int_{0}^{+\infty} e^{-a x^2} \mathrm{d}x = \frac{1}{2}\sqrt{\frac{\pi}{a}} \quad (a > 0)
            \]

            \item \textbf{含 $x$ 的推广形式}：
            \begin{align*}
                \int_{0}^{+\infty} x e^{-a x^2} \mathrm{d}x &= \frac{1}{2a} \quad (a > 0) \\
                \int_{0}^{+\infty} x^2 e^{-a x^2} \mathrm{d}x &= \frac{1}{4}\sqrt{\frac{\pi}{a^3}} \quad (a > 0) \\
                \int_{-\infty}^{+\infty} x^2 e^{-a x^2} \mathrm{d}x &= \frac{1}{2a}\sqrt{\frac{\pi}{a}} \quad (a > 0)
            \end{align*}

            \item \textbf{应用}：高斯积分是概率统计中正态分布（高斯分布）的基础，在物理学、工程学等领域有广泛应用
        \end{circlenum}
    \end{enumerate}

    \subsection{公式}

    \begin{enumerate}
        \item $\displaystyle\int_0^x \sin \frac{y}{t}\,\mathrm{d}y = \int_0^x t\cos \frac{y}{t}\,\mathrm{d}y$
        \item $\displaystyle\int_0^1 \arcsin \sqrt {1-x^2}\mathrm{d}x = 1$
        \item $\displaystyle\int_0^1 x\arcsin \sqrt {4-4x^2}\mathrm{d}x = \frac{1}{2}$
        \item $\displaystyle\int \frac{1}{3\cos^2 \theta + \sin^2 \theta} = \int \frac{\mathrm{d}(\tan \theta)}{3 + \tan^2 \theta} = \frac{1}{\sqrt {3}}\arctan\frac{\tan \theta}{\sqrt {3}} + C$
        \item $\displaystyle\int e^{-x}\mathrm{d}x = -e^{-x} + C$
        \item $\Gamma$函数：
        \[
            \begin{aligned}
                &\Gamma(\alpha) = \int_{0}^{+\infty} x^{\alpha-1} e^{-x}\,\mathrm{d}x = 2\int_{0}^{+\infty} t^{2\alpha-1} e^{-t^2}\,\mathrm{d}t (\alpha, t > 0) \\
                &\Gamma(\alpha+1) = \alpha \Gamma(\alpha) \\
                &\Gamma(1) = 1 \quad \Gamma\!\left(\frac{1}{2}\right) = \sqrt{\pi}
            \end{aligned}
        \]
        \item $\displaystyle(\frac{\sin \theta}{\cos \theta + \sin \theta})' = \frac{1}{(\cos \theta + \sin \theta)^2}$
    \end{enumerate}

    \subsection{方法总结}

    \begin{enumerate}
        \item 寻找二重积分最大值的区域，即是寻找被积函数大于等于零的区域
        \item 研究某一点处的导数，需使用导数的定义
        \item 当 \begin{cases}
                     \text{被积函数被命制成具体函数，但}\iint\limits_D f(x,y) \mathrm{d}\sigma\text{难计算 } \\
                     \text{被积函数被命制成抽象函数}
        \end{cases}时，可以利用二重积分的中值定理来处理
        \item 遇到二重积分的题目，首先考虑对称性，看能否化简
        \item 当题目是累次积分时，通常考虑换积分顺序或者坐标系
        \item 交换积分顺序的题目，应先确定积分区域（注意{\color{red}{上下限大小}}），然后交换积分顺序写成相应的累次积分
        \[
            \begin{aligned}
                &\int_{\frac{\pi}{2}}^{\pi} \mathrm{d}x \int_{1}^{\sin x} f(x, y)\,\mathrm{d}y \quad \text{由于 } \sin x < 1,\ \text{该式不是标准的二重积分，仅为累次积分形式} \\
                &= - \int_{\frac{\pi}{2}}^{\pi} \mathrm{d}x \int_{\sin x}^{1} f(x, y)\,\mathrm{d}y \quad \text{此时为对应区域上的二重积分}
            \end{aligned}
        \]
        \item $\displaystyle \frac{\sin x}{x}, \frac{\cos x}{x}, \frac{\ln (1 + x)}{x}, \frac{1}{\ln x}, \sin x^2, \cos x^2, \sin \frac{1}{x}, \cos \frac{1}{x}, \frac{\tan x}{x}, \frac{e^x}{x}, \tan x^2, e^{ax^2 + bx + c}(a \neq 0) $均没有初等函数形式的原函数，见到它们，一般都要{\color{red}{交换积分次序}}，避免先对$x$求积分
        \item 积分与字母用谁无关
        \begin{itemize}
            \item $\displaystyle\iint_{D_{xy}} f(x, y)\,\mathrm{d}x\mathrm{d}y = \iint_{D_{yx}} f(y, x)\,\mathrm{d}y\mathrm{d}x$
        \end{itemize}
        \item 当分母$\to 0$时，若分式的极限不为$0$，则可知分子$\to 0$
        \item 思考：$\displaystyle\lim\limits_{x\to +\infty} e^{x^2}(\int_0^x t^2e^{-t^2}\,\mathrm{d}t + a) = \infty$
    \end{enumerate}

    \subsection{理解}

    \begin{enumerate}
        \item 设平面区域$D = \{(x, y)|0\leq x \leq 1-y, 0 \leq y \leq 1\}$，计算二重积分$\displaystyle\iint_D e^{\frac{y}{x+y}}\,\mathrm{d}\sigma$
        \begin{analysisbox}
            \textbf{分析：}
            区域 $D$ 由 $x\ge 0,\ y\ge 0,\ x+y\le 1$ 给出，是第一象限内三角形。按
            \[
                u=x+y,\qquad v=y
            \]
            换元后，被积函数化为 $e^{v/u}$，区域变为 $0\le v\le u\le 1$，积分可直接计算。

            \textbf{解：}令\begin{cases}
                              x + y = u, \\
                              y = v
            \end{cases} 则 \begin{cases}
                               x = u - v, \\
                               y = v
            \end{cases}
            \[
                J=\frac{\partial(x,y)}{\partial(u,v)}
                =\begin{vmatrix}
                     1 & -1 \\
                     0 & 1
                \end{vmatrix}
                =1.
            \]
            由区域条件 $x\ge0,\ y\ge0,\ x+y\le1$ 可得$0\le v\le1,\qquad v\le u\le1$，因而
            \[
                \begin{aligned}
                    \iint\limits_D e^{\frac{y}{x+y}}\,\mathrm{d}\sigma
                    &= \iint\limits_{D_{uv}} e^{\frac{v}{u}} |J|\,\mathrm{d}v\,\mathrm{d}u
                    = \int_0^1 \mathrm{d}u \int_0^u e^{\frac{v}{u}} \,\mathrm{d}v \\
                    &= \int_0^1 \left. e^{\frac{v}{u}} \right|_{v=0}^{v=u} \mathrm{d}u
                    = \int_0^1 u(e-1)\,\mathrm{d}u
                    = \frac{1}{2}(e-1)
                \end{aligned}
            \]
        \end{analysisbox}
        \item 已知$\lim\limits_{x\to +\infty}\displaystyle\frac{\int_0^x t^2e^{x^2 -t^2}\,\mathrm{d}t + ae^{x^2}}{x^b} = -\frac{1}{2}$，求$a, b$的值
        \begin{analysisbox}
            \textbf{分析：}
            将分子提取公因子 $e^{x^2}$，先判断常数项是否会导致指数爆炸；若要极限为有限常数，必须先消去主导项。随后把差值改写为尾积分 $\int_x^{+\infty} t^2e^{-t^2}\,\mathrm{d}t$，再做渐近估计确定幂次 $b$。

            \textbf{解：}
            设
            \[
                I(x)=\int_0^x t^2e^{-t^2}\,\mathrm{d}t,\qquad I_\infty=\int_0^{+\infty} t^2e^{-t^2}\,\mathrm{d}t.
            \]
            原式可写为
            \[
                \frac{\int_0^x t^2e^{x^2-t^2}\,\mathrm{d}t+ae^{x^2}}{x^b}
                =\frac{e^{x^2}\bigl(I(x)+a\bigr)}{x^b}.
            \]
            当 $x\to+\infty$ 时，$I(x)\to I_\infty$，且
            \[
                I_\infty=\int_0^{+\infty} t^2e^{-t^2}\,\mathrm{d}t
                =\frac{\sqrt{\pi}}{4}.
            \]
            若 $a+I_\infty\neq0$，则分子与 $e^{x^2}$ 同阶，除以任意 $x^b$ 都不可能收敛到有限常数，矛盾。
            故必须有
            \[
                a=-I_\infty=-\frac{\sqrt{\pi}}{4}.
            \]
            此时
            \[
                e^{x^2}\bigl(I(x)+a\bigr)
                =e^{x^2}\bigl(I(x)-I_\infty\bigr)
                =-e^{x^2}\int_x^{+\infty} t^2e^{-t^2}\,\mathrm{d}t.
            \]
            对尾积分分部积分：
            \[
                \int_x^{+\infty} t^2e^{-t^2}\,\mathrm{d}t
                =\frac{x}{2}e^{-x^2}+\frac12\int_x^{+\infty}e^{-t^2}\,\mathrm{d}t
                =\frac{x}{2}e^{-x^2}+o\!\bigl(e^{-x^2}x\bigr).
            \]
            因而
            \[
                e^{x^2}\bigl(I(x)+a\bigr)
                =-\frac{x}{2}+o(x).
            \]
            所以原极限为
            \[
                \lim_{x\to+\infty}\frac{-\frac{x}{2}+o(x)}{x^b}=-\frac12.
            \]
            由此只能取 $b=1$。

            综上，
            \[
                \boxed{a=-\frac{\sqrt{\pi}}{4},\quad b=1.}
            \]
        \end{analysisbox}
    \end{enumerate}

    \subsection{三换（数学核心方法论）}
    数学问题的本质往往不变，但表现形式可以转换。「三换」是指通过{\color{red}{改变问题的表达形式}}来简化求解过程的三种基本策略。

    \begin{circlenum}
        \item \textbf{换序}：交换运算次序

            {\color{blue}核心思想}：当一种运算次序难以进行时，尝试交换次序

            {\color{blue}典型应用}：
        \begin{itemize}
            \item 积分换序：$\displaystyle\int\mathrm{d}x\int\mathrm{d}y \leftrightarrow \int\mathrm{d}y\int\mathrm{d}x$
            \item 求和与极限交换：$\displaystyle\lim_{n\to\infty}\sum_{k=1}^{n} \leftrightarrow \sum_{k=1}^{\infty}\lim_{n\to\infty}$
            \item 微分与积分交换：$\displaystyle\frac{\mathrm{d}}{\mathrm{d}x}\int \leftrightarrow \int\frac{\mathrm{d}}{\mathrm{d}x}$
        \end{itemize}

        \item \textbf{换系}：转换坐标系/表达系统

            {\color{blue}核心思想}：同一数学对象在不同系统中有不同表达，选择最简形式

            {\color{blue}典型应用}：
        \begin{itemize}
            \item 坐标系变换：直角坐标 $\leftrightarrow$ 极坐标 $\leftrightarrow$ 球坐标
            \item 域变换：时域 $\xleftrightarrow{\text{傅里叶}}$ 频域
            \item 基变换：线性空间中不同基下的坐标转换
        \end{itemize}

        \item \textbf{换元}：变量替换

            {\color{blue}核心思想}：用新变量替代原变量，将复杂形式化为简单形式

            {\color{blue}典型应用}：
        \begin{itemize}
            \item 不定积分换元：$\displaystyle\int f[\varphi(x)]\varphi'(x)\mathrm{d}x = \int f(u)\mathrm{d}u$
            \item 二重积分换元：雅可比行列式变换
            \item 方程化简：通过适当换元将方程化为标准形式
        \end{itemize}
    \end{circlenum}

    \vspace{0.5em}
    {\color{red}$\bigstar$} \textbf{三换的本质}：问题本身不变，改变的是我们{\color{red}{观察问题的角度}}和{\color{red}{处理问题的工具}}。选择合适的形式，往往能让困难问题迎刃而解。

\end{document}
